\documentclass[a4paper,10pt]{report}\input{../0.Préambule/Préambule_cours_prof}\begin{document}

\pagestyle{empty}
\newpage
\subsection*{La pyramide de Khéops}
\addcontentsline{toc}{section}{La pyramide de Khéops}

\noindent \begin{minipage}{8cm}
La pyramide de Khéops a été construite il y a plus de $\np{4 500}$ ans en Égypte. C'est la plus grande des fameuses pyramides de Gizeh, près du Caire. 

\medskip
Sa base est carrée et mesure $230$ mètres. Sa hauteur est de $137$ mètres. On a calculé  qu'il avait fallu $\np{2 300 000}$ blocs de pierre d'une masse moyenne de $2,3$t pour l'édifier.

\medskip 
Napoléon a même calculé qu'avec ces blocs, on aurait pu édifier un  mur de $3$ mètres de haut et de $30$cm d'épaisseur le long des frontières de la France. 
\end{minipage}
\begin{minipage}{9cm}
\begin{center}
\includegraphics[scale=0.38]{../40.Image/08. Act - Espace - Aires et volumes 1.png}
\end{center}
\end{minipage}

\medskip
Il aura fallu $20$ années et une main d'œuvre de $\np{10 000}$ hommes renouvelés tous les trois mois selon Hérodote pour construire cet édifice. 

\begin{enumerate}
\item Combien d'hommes auraient participé à la construction de la pyramide ?
\item Quelle est l'aire de la base de la pyramide ?
\item Quel est le volume de la pyramide ?
\item Quel est le volume moyen d'un bloc de pierre ?
\item Quelle est la masse de la pyramide ?
\item La taille des blocs utilisés à la base de la pyramide mesurent en moyenne 2,50 m de
long, 1 m de large et 1,20 m de haut.
\begin{enumerate}
\item Combien de blocs a-t-il fallu utiliser pour construire la base de
la pyramide ?
\item Quel volume ces pierres représentent-elles ?
\end{enumerate}
\end{enumerate}

\subsection*{La pyramide de Khéops}
\addcontentsline{toc}{section}{La pyramide de Khéops}

\noindent \begin{minipage}{8cm}
La pyramide de Khéops a été construite il y a plus de $\np{4 500}$ ans en Égypte. C'est la plus grande des fameuses pyramides de Gizeh, près du Caire. 

\medskip
Sa base est carrée et mesure $230$ mètres. Sa hauteur est de $137$ mètres. On a calculé  qu'il avait fallu $\np{2 300 000}$ blocs de pierre d'une masse moyenne de $2,3$t pour l'édifier.

\medskip 
Napoléon a même calculé qu'avec ces blocs, on aurait pu édifier un  mur de $3$ mètres de haut et de $30$cm d'épaisseur le long des frontières de la France. 
\end{minipage}
\begin{minipage}{9cm}
\begin{center}
\includegraphics[scale=0.38]{../40.Image/08. Act - Espace - Aires et volumes 1.png}
\end{center}
\end{minipage}

\medskip
Il aura fallu $20$ années et une main d'œuvre de $\np{10 000}$ hommes renouvelés tous les trois mois selon Hérodote pour construire cet édifice. 

\begin{enumerate}
\item Combien d'hommes auraient participé à la construction de la pyramide ?
\item Quelle est l'aire de la base de la pyramide ?
\item Quel est le volume de la pyramide ?
\item Quel est le volume moyen d'un bloc de pierre ?
\item Quelle est la masse de la pyramide ?
\item La taille des blocs utilisés à la base de la pyramide mesurent en moyenne 2,50 m de
long, 1 m de large et 1,20 m de haut.
\begin{enumerate}
\item Combien de blocs a-t-il fallu utiliser pour construire la base de
la pyramide ?
\item Quel volume ces pierres représentent-elles ?
\end{enumerate}
\end{enumerate}
\end{document}